% 第三章按用户旧版结构修订。模型采用auto_low_immediate_v1，实验使用既有已归档结果。

\chapter[多场景端侧推理的热约束节能联合规划]{多场景端侧推理的热约束\\节能联合规划}

\label{chap:model}



\section{本章引言}



端侧设备可以利用CPU、GPU和NPU协同执行多个推理应用，不同应用组合对计算资源的需求也有所不同。例如，设备在一种场景下同时运行应用A与B，在另一种场景下运行应用A与C。应用A虽然执行相同的计算阶段，但与其竞争资源的应用发生了变化，适合的频率和执行安排因而可能不同。为使部署适用于多个场景，本章采用跨场景共享Stage映射的设计方式：同一应用的固定计算阶段绑定同一执行单元，各场景分别选择工作点和执行时序。Kim等采用共同映射与逐场景调频组织多应用部署\cite{kim2023scenario}，为本章的场景关联方式提供了参考。



共享部署的节能效果需要结合各场景的完整执行过程评价。一个阶段单独使用某设备时能量较低，并不意味着多个请求共用该设备后仍能按时完成；降低频率也可能延长资源占用，而且不同阶段的时长增幅可以不同。与此同时，功率在时间和空间上的分布决定温度变化，温度又通过漏电影响功耗。已有研究分别讨论了任务相关频率响应、映射调频与热状态之间的联系\cite{source2,source1,fdtm2022}。在多场景共享部署条件下，需要将这些关系放到同一执行过程中，判断部署能否同时适应不同场景的资源竞争、请求期限和入口温度。



基于上述关系，本章研究多场景端侧推理的热约束节能联合规划问题。给定各场景的应用组合和有限请求，在规划阶段共同选择Stage的执行单元，并分别确定各场景的工作点与启动时刻，使全部场景满足期限及连续温度约束，最小化按给定权重汇总的系统能耗。模型保留逐Stage的频率响应、共享频率域、域空闲立即采用最低工作点以及温度相关功耗，使能量评价包含工作、等待和完成后的空闲过程。



本章首先介绍任务与场景、处理器及执行模式，随后建立功耗、能耗和瞬态温度模型。在此基础上定义联合优化问题，并针对有限启动集合给出MILP表达及连续热反馈求解流程。最后结合已经完成的仿真结果，分析响应参数、完整窗口能耗和温度安全对调度安排的影响。



\section{系统模型}



多场景联合规划需要先明确各场景执行哪些工作、哪些部署需要共用，以及每项计算有哪些合法执行选择。本节依次建立任务与场景模型、处理器模型和执行模式模型。



\subsection{任务与多场景建模}



设推理应用集合为\(\mathcal A\)。每个应用的模型、输入规格、精度和Stage划分在规划前确定，其固定工作流表示为有向无环图

\begin{equation}
G_a=(V_a,E_a),\qquad a\in\mathcal A,
\label{eq:application-dag}
\end{equation}

其中，\(V_a\)为Stage集合，\(E_a\)为数据依赖边集合。一个Stage包含在同一执行单元上连续完成的一组算子；若\((i,j)\in E_a\)，则后继Stage需要等待前驱结果。该表示允许链式、分支及汇合结构，不预先限定Stage数量。



端侧服务需求变化时，同时运行的应用组合以及各应用的请求数量、到达和期限可能不同。例如，一种场景同时运行应用A与B，另一种场景运行A与C。每个应用在一个场景中可以被调用多次，每次调用构成一个请求。若应用A包含\(A_1\rightarrow A_2\)，两次请求就分别产生一套\(A_1\)、\(A_2\)的执行实例；它们来自相同工作流，但分别计算并占用资源。



为描述这些已知需求，设有限场景集合为\(\mathcal C\)，场景\(c\)包含请求集合\(\mathcal R_c\)。以\(a(r)\)表示请求所属应用，\(r_{cr}\ge0\)表示释放时刻，\(d^{\rm rel}_{cr}>0\)表示相对截止时间，则请求描述及绝对截止时间为

\begin{equation}
\bigl(a(r),r_{cr},d^{\rm rel}_{cr}\bigr),\qquad
D_{cr}=r_{cr}+d^{\rm rel}_{cr}.
\label{eq:request-description}
\end{equation}

本章在执行前获得全部场景及请求信息。场景改变的是应用共存关系和服务条件，应用本身的固定工作流保持不变。



每个请求展开为独立的Stage实例图。记场景\(c\)的全部实例为\(\mathcal I_c\)，依赖边为\(\mathcal E_c\)，以\(r(c,i)\)表示实例\(i\)所属请求，以\(\mathcal O_{cr}\)表示请求必须产生的输出实例集合。不同请求即使调用同一应用，也分别执行并占用资源；请求之间仅在存在实际数据依赖时连接依赖边。请求完成意味着\(\mathcal O_{cr}\)中的全部输出均已完成，不能仅以其中一个分支结束代替完整请求完成。



为关联不同场景中的部署，定义固定Stage身份集合

\begin{equation}
\mathcal J=\{(a,j):a\in\mathcal A,\ j\in V_a\},\qquad
\tau(c,i)\in\mathcal J.
\label{eq:stage-identity}
\end{equation}

\(\tau(c,i)\)指出实例来自哪个应用的哪个Stage。同一身份\(\tau\)在所有场景和请求中使用共同执行单元\(u_\tau\)，工作点与启动时刻则可以分别选择。在上述例子中，若将\(A_1\)部署到NPU0，则A＋B和A＋C中每次请求的\(A_1\)都使用NPU0，具体工作点和起点再根据各场景的竞争与期限确定。因此，共享部署关联的是同一Stage的执行位置，各次请求仍独立完成计算。各场景从给定入口状态分别评价，场景之间通过共同部署建立联系。



\subsection{处理器模型}



平台包含CPU、GPU和NPU三类执行资源，分别记为\(\mathcal U_C\)、\(\mathcal U_G\)和\(\mathcal U_N\)，全部执行单元为\(\mathcal U\)。CPU单元可为单核或固定核组，各单元占用的核集合互不重叠，线程数和亲和性预先确定。GPU与NPU按设备映射，其内部并行能力体现在Stage响应中。每个执行单元同时至多执行一个Stage，不同单元可以并发，设备数量与合法支持范围由平台输入给出。



处理器通过合法的电压--频率工作点运行。设频率控制域集合为\(\mathcal H\)，\(\gamma(u)\)为单元\(u\)所属的域，\(\mathcal K_h\)为域\(h\)的有限工作点集合。工作点\(k\)给出配套的频率\(f_h(k)\)与电压\(v_h(k)\)，二者不独立选择。同域CPU单元共用控制状态，具有独立控制权限的GPU或NPU使用相应独立域。仅能固定运行的设备可用单元素工作点集合表示。



执行单元决定Stage占用哪份计算资源，频率域决定哪些单元必须共同调频。例如，CPU1与CPU2可以分别执行两个Stage，但若属于同一频率域，并发期间就必须采用同一工作点；CPU1先完成时，只要CPU2仍在执行，该域就不能换点。



Stage执行期间不抢占、不迁移，也不改变工作点。将一批工作启动至该域工作全部完成的执行过程称为忙段，期间加入的工作也需完成。同域内同时执行的Stage必须采用同一个相容域状态；只要仍有Stage持续执行，该域就不能换点。域内旧工作全部完成后，可以在下一批工作启动时选择新的工作点，两批之间不必插入正长度空闲。因此，各忙段可以采用不同工作点，同一Stage的执行过程则保持工作点不变。



空闲频率控制与忙段中的执行状态相衔接。设$L_h$为频率域$h$唯一的最低工作点，$k_{ch}(t)$为域状态轨迹。本章采用空闲自动最低档（auto-low）策略：域中只要没有正在执行的Stage，就立即切换至$L_h$；再次启动Stage时按照选中模式选择工作点。同域多个Stage并发时仍使用相同工作点，任何持续执行的Stage均不会在运行中换档。各场景入口没有未完成工作，控制域起始于最低档；最后一次任务完成后，直至窗口结束的空闲尾段也保持在最低档。这里采用零额外搬运、交接和换档时间及能量开销的模型假设，应用工作流本身包含的预处理与后处理计算仍完整保留。

\subsection{执行模式与时长模型}



执行单元与合法工作点共同决定Stage的运行方式。将二者的组合称为执行模式，场景\(c\)中实例\(i\)的合法模式集合记为\(\mathcal M_{ci}\)。模式\(m\)指定单元\(u(m)\)和完整域工作点\(\kappa(m)\)；对应时长由响应参数计算，活动功率由匹配功耗参数给出。例如，同一Stage可以在CPU1的不同工作点或某个支持它的NPU工作点上执行，每种选择具有各自的时间与功率代价。规划器选择模式，据此采用相应参数。只有后端支持该Stage且控制要求相容的组合才进入合法集合。



不同Stage对降频的敏感程度并不相同。借鉴可缩放部分与固定部分的时长分解\cite{source2}，对固定Stage \(\tau\)在单元\(u\)上的响应采用

\begin{equation}
d_{\tau u}(f)=d^{\rm ref}_{\tau u}
\left[(1-\alpha_{\tau u})+\alpha_{\tau u}\frac{f_u^{\rm ref}}{f}\right],
\qquad 0\le\alpha_{\tau u}\le1.
\label{eq:response-model}
\end{equation}

其中，\(d^{\rm ref}_{\tau u}\)为参考频率\(f_u^{\rm ref}\)下的时长，\(\alpha_{\tau u}\)为参考时长中随频率反比变化的比例。两者均为Stage与设备相关的输入参数，在声明的频率范围内使用。\(\alpha=1\)对应反比缩放，\(\alpha=0\)对应时长不随频率改变；频率减半时，时长变为参考值的\(1+\alpha\)倍。因此，相同降频比例可以对不同Stage产生不同的资源占用代价。



对模式\(m\)，将其合法频率代入式\eqref{eq:response-model}即可得到\(d_\tau(m)>0\)。活动增量功率\(q_\tau(m)\ge0\)另由匹配功耗参数给出，不能从时长公式推算。相同Stage、工作量、单元和工作点跨场景复用相同参数；当前模型中，这些时长和活动增量不随并发组合或温度变化。参考点及响应参数须配套确定，改变参考频率时不能只替换公式中的频率而保持其余参数不变。



记实例选中的模式为\(m_{ci}\)，启动时刻为\(s_{ci}\)。其结束时刻、半开执行区间和活动指示分别为

\begin{equation}
\begin{aligned}
F_{ci}&=s_{ci}+d_{\tau(c,i)}(m_{ci}),\\
I_{ci}&=[s_{ci},F_{ci}),\qquad
\iota_{ci}(t)=\mathbf1_{I_{ci}}(t).
\end{aligned}
\label{eq:execution-interval}
\end{equation}

半开区间允许一个Stage结束时另一个立即开始。模式与起点确定后，结合给定入口和空闲自动最低档规则，即可导出执行次序、并发、等待及域工作点轨迹，构成后续能耗与温度评价的时间轴。



\section{功耗与温度模型}



设备工作时产生热量，温度变化又影响漏电功耗。为计算这一相互作用，本节先说明工作与空闲的功率组成，再分别给出三类设备的计量方式，最后将完整功率轨迹与热网络连接。



\subsection{功耗组成与运行状态}



处理器功耗通常包括动态功耗与静态漏电功耗，其基本关系可写为\cite{source1}

\begin{equation}
P_{\rm dyn}\approx aC^{\rm eff}v^2f,\qquad
P_{\rm leak}=vI_{\rm leak}(v,T).
\label{eq:physical-power}
\end{equation}

其中，\(a\)为开关活动因子，\(C^{\rm eff}\)为等效电容，\(v\)和\(f\)为电压与频率。工作点改变动态功耗和漏电水平，温度进一步改变漏电流。温度相关功率与热状态的联合描述已用于多核及CPU--GPU热管理\cite{chantem2011,fdtm2022}。



动态功耗与漏电功耗解释功率的物理来源。为计算Stage执行及其间隙中的耗电，本章进一步按运行状态将功率组织为参考基准、温度修正和活动增量。参考基准是在给定工作点和参考温度下持续存在的功率，已包含该条件下的漏电；温度修正补充当前温度相对参考值引起的变化；活动增量只在Stage执行时加入。空闲时仍可能存在开关活动，工作时也仍有漏电，因此运行状态与物理功耗来源应分别描述。



以独占一个计量组、对应一个热节点的设备为例。在固定工作点下，以\(B\)表示参考基准、\(\lambda\)表示温度系数、\(q\)表示所执行Stage的活动增量，则空闲功率为\(B+\lambda(T-T^{\rm ref})\)，工作功率为\(B+\lambda(T-T^{\rm ref})+q\)。Stage结束时活动项消失，基准和随温度变化的部分仍持续存在。下面将这一关系推广到包含共享基准和多个热节点的计量边界。



设热节点集合为\(\mathcal N\)，互不重叠的功率计量组为\(\mathcal B\)，\(\pi(u)\)给出执行单元所属组，\(\xi_{cg}(t)\)表示影响组\(g\)功耗的全部域状态组合。组参考基准记为\(B_g(\xi_{cg})\)，参考温度为\(T_g^{\rm ref}\)。组内位于节点\(n\)的温度修正为

\begin{equation}
\Delta P^\lambda_{cgn}(t)=
\lambda_{gn}(\xi_{cg}(t))\,[T_{cn}(t)-T_g^{\rm ref}],
\qquad\lambda_{gn}\ge0.
\label{eq:local-leakage}
\end{equation}

\(\lambda_{gn}\)的单位为W/K，使用该来源所在节点的本地温度；不存在对应来源时取零。系数随硬件与工作点给定，工作和空闲使用同一系数。该线性关系在声明有效温区内使用。低于参考温度时修正可为负，表示漏电较参考值减小，完整组功率与节点功率仍须满足物理有效性要求。



\subsection{CPU功耗模型}



多个CPU执行单元可以并行工作，因此需要把各单元的活动功率相加，同时避免重复计入共享基准。令\(u_{ci}=u(m_{ci})\)、\(q_{ci}=q_{\tau(c,i)}(m_{ci})\)。计量组\(g\)中的CPU活动增量为

\begin{equation}
\Delta P^{C}_{cg}(t)=
\sum_{\substack{i\in\mathcal I_c:\,u_{ci}\in\mathcal U_C,\ \pi(u_{ci})=g}}
q_{ci}\iota_{ci}(t).
\label{eq:cpu-active-power}
\end{equation}

组基准\(B_g\)只计一次，而各CPU单元的温度修正分别使用自身节点温度。共用频率域限制的是工作点选择，并不要求两个单元温度相同。活动增量可以由同工作点、同温度和同一计量边界下的工作功率与空闲功率之差获得，共享组内其他活动保持一致。这里扣除的空闲功率包含该温度下的修正项，不能仅扣参考温度下的\(B_g\)。也可由具有明确适用范围的功耗模型直接提供活动增量。



\subsection{GPU功耗模型}



GPU以设备作为可调度单元，其内部计算与存储活动汇总到Stage活动增量中。对设备\(u\in\mathcal U_G\)，有

\begin{equation}
\Delta P^{G}_{cu}(t)=
\sum_{i\in\mathcal I_c:\,u_{ci}=u}q_{ci}\iota_{ci}(t).
\label{eq:gpu-active-power}
\end{equation}

设备独占规则使上式在同一时刻至多包含一个活动Stage。GPU基准随当前工作点确定，温度修正使用设备节点温度。F-DTM按设备和频率给出温度相关功率\cite{fdtm2022}，本章沿用这一物理联系，将Stage活动增量与持续基准分开计量。



\subsection{NPU功耗模型}



NPU同样采用设备级计量，区别主要在合法后端支持和设备响应。Stage只有在后端能够完整执行其算子范围时才可使用该NPU；内部并行能力和算子活动体现于时长及功率参数。对\(u\in\mathcal U_N\)，活动增量为

\begin{equation}
\Delta P^{N}_{cu}(t)=
\sum_{i\in\mathcal I_c:\,u_{ci}=u}q_{ci}\iota_{ci}(t).
\label{eq:npu-active-power}
\end{equation}

NPU可从其合法工作点中选择，基准和温度系数使用设备自身参数。若一台加速器独占一个计量组，设备功率就是该组基准、本地温度修正与式\eqref{eq:gpu-active-power}或式\eqref{eq:npu-active-power}之和；若计量边界包含共享部分，则按下一节的组表达汇总。三类处理器使用相同物理分解，但不共用参数，也不需要人为设置不同的功耗定律。



\subsection{系统功耗与能耗}



\textbf{（1）完整能耗计算。}

将共享基准、本地温度修正与实际映射的活动增量汇总，计量组电功率为

\begin{equation}
P_{cg}(t)=B_g(\xi_{cg}(t))
+\sum_{n\in\mathcal N}\Delta P^\lambda_{cgn}(t)
+\sum_{i\in\mathcal I_c:\pi(u_{ci})=g}q_{ci}\iota_{ci}(t),
\label{eq:group-electric-power}
\end{equation}

其中，\(\mathcal N\)为热节点集合。计量组用于确定基准计费边界，频率域用于确定共同控制关系，热节点用于记录空间温度。例如，两个CPU单元可以共用一个频率域和一份组基准，但各自对应不同热节点：工作点共同选择，基准只计一次，温度分别计算。三者的对应关系由平台输入给出。



不同安排可能在不同时刻完成，但能耗比较需要采用相同的时间边界。因此，为场景\(c\)给定固定评价窗口\([0,W_c]\)，满足\(W_c\ge\max_{r\in\mathcal R_c}D_{cr}\)。窗口长度由工作负载与评价需求给定，不随候选完成时间改变。设\(P_{{\rm ext},c}(t)\)为计入系统能耗、位于当前热网络边界外的给定公共功率，则

\begin{equation}
\begin{aligned}
P_{{\rm sys},c}(t)&=P_{{\rm ext},c}(t)+\sum_{g\in\mathcal B}P_{cg}(t),\\
E_c(m_c,s_c)&=\int_0^{W_c}P_{{\rm sys},c}(t)\,\mathrm dt,\\
\bar E&=\sum_{c\in\mathcal C}p_cE_c(m_c,s_c),\quad
p_c\ge0,\quad\sum_cp_c=1.
\end{aligned}
\label{eq:full-energy}
\end{equation}

\(p_c\)为预先给定的完整场景片段发生次数权重，所有候选采用同一组权重。等待及请求完成后的尾段仍有基准和温度相关功耗，均计入积分。模式固定时，活动能量\(\sum_iq_{ci}d_{\tau(c,i)}(m_{ci})\)已经确定；调整起点却可能改变共享频率域的忙段长度、较高档位的驻留时间和温度轨迹，因此完整能量一般仍依赖时序。



\textbf{（2）功率到热节点的对应。}

完整能耗汇总功率的总量，热计算还需要确定这些功率在哪里产生。记\(n(u)\)为执行单元对应节点，活动注热映射取\(h_{nu}=\mathbf1_{n=n(u)}\)。共享基准采用输入给定的分布\(a^B_{ng}(\xi_g)\ge0\)，对边界内组满足\(\sum_na^B_{ng}=1\)。节点注热为

\begin{equation}
p^{\rm heat}_{cn}(t)=
\sum_g a^B_{ng}(\xi_{cg}(t))B_g(\xi_{cg}(t))
+\sum_g\Delta P^\lambda_{cgn}(t)
+\sum_i h_{n,u_{ci}}q_{ci}\iota_{ci}(t).
\label{eq:local-heat-input}
\end{equation}

温度修正注入其来源节点，基准另按给定分布处理，不默认平均分配。由此有\(\sum_np^{\rm heat}_{cn}=\sum_gP_{cg}\)，热分配不产生第二份能耗。边界外功率仅进入式\eqref{eq:full-energy}；节点间传热由热网络描述，不通过分摊活动功率代替。



\subsection{瞬态温度模型}



热节点的温度既取决于当前注热，也取决于入口储热和节点间导热。CPU按固定执行单元区分节点，GPU与NPU按设备区分节点。采用热阻热容网络描述这些关系\cite{chantem2011,fdtm2022}。设环境温度\(T_a\)恒定，温升向量为\(\theta_c=T_c-T_a\mathbf1\)，热容矩阵\(C=\operatorname{diag}(C_n)\)满足\(C_n>0\)。以\(g_{nj}=g_{jn}\ge0\)表示节点间导热，以\(g_{na}\ge0\)表示对环境导热，定义

\begin{equation}
G_{nn}=g_{na}+\sum_{j\ne n}g_{nj},\qquad
G_{nj}=-g_{nj}\ (j\ne n).
\label{eq:conductance-matrix}
\end{equation}

该表达描述被动导热网络。热容单位为J/K，导热单位为W/K。由能量平衡得到

\begin{equation}
C\dot\theta_c(t)=-G\theta_c(t)+p^{\rm heat}_c(t),\qquad
\theta_c(0)=T_c^0-T_a\mathbf1.
\label{eq:thermal-ode}
\end{equation}

\(T_c^0\)为场景给定入口温度。由于式\eqref{eq:local-heat-input}含本地温度修正，注热与温度相互作用；工作点和活动状态固定的片段内，方程为常系数仿射系统。工作点或活动集合改变时更新片段系数，温度在边界连续。被动导热网络本身的稳定性不代替含漏电反馈后的温度检查。



每个节点给定温度上限\(T_n^{\max}\)。完整安排必须满足

\begin{equation}
T_a+\theta_{cn}(t)\le T_n^{\max},\qquad
n\in\mathcal N,\quad t\in[0,W_c].
\label{eq:continuous-temperature-limit}
\end{equation}

功耗与热参数还给定模型有效范围，温度轨迹和完整物理功率须始终处于该范围内。检查覆盖入口、工作、等待和空闲尾段，也包括事件之间的温度峰值。时间有限的窗口保证不自动延伸为窗口之后的持续安全。



\section{问题定义}



前两节给出了工作需求、执行规则和热能耗关系，本节据此明确共同优化的对象。输入按用途分为表\ref{tab:problem-inputs}中的五类，均在规划前给定。规划器选择部署和执行安排，响应参数、物理参数及评价条件保持为输入。



\begin{table}[htbp]
\centering
\caption{联合规划问题的输入}
\label{tab:problem-inputs}
\begin{tabular}{p{0.17\textwidth}p{0.74\textwidth}}
\toprule
类别 & 主要内容 \\
\midrule
工作需求 & 各场景请求DAG、释放时刻、逐请求期限及固定Stage身份。\\
执行选择 & 各Stage的合法执行单元与工作点组合。\\
平台规则 & 执行资源、频率域及其控制规则。\\
性能与物理参数 & 参考时长、响应系数、功耗参数和热网络参数。\\
评价条件 & 各场景入口温度、评价窗口、权重及空闲最低档策略。\\
\bottomrule
\end{tabular}
\end{table}



共同决策按作用范围分为两类。跨场景的\(u=(u_\tau)_{\tau\in\mathcal J}\)确定各固定Stage共享的具体执行单元；场景内的\(m_c=(m_{ci})\)与\(s_c=(s_{ci})\)选择各实例的执行模式和启动时刻。模式包含设备与工作点，其中设备分量须满足\(u(m_{ci})=u_{\tau(c,i)}\)，工作点则可结合该场景分别选择。



这两类决定在同一优化问题中共同确定，不规定先固定部署再规划时序。例如，\(A_1\)在A＋B与A＋C中使用相同设备，但可采用不同工作点和起点；共享设备选择需要结合两种场景的安排一起评价。结束时刻、控制轨迹、温度和能量由这些选择导出，后续即使以辅助变量表示，也必须满足相应执行与物理关系。



以\(\mathcal F\)表示满足合法模式、共享部署、请求释放与依赖、执行单元互斥、工作点控制、逐请求期限、窗口覆盖、连续热约束及物理有效条件的完整安排集合。一般问题定义为

\begin{equation}
\begin{aligned}
(P)\qquad \min_{u,m,s}\quad &\sum_{c\in\mathcal C}p_cE_c(m_c,s_c)\\
\mathrm{s.t.}\quad &(u,m,s)\in\mathcal F,\qquad
s_{ci}\in\mathbb R_{\ge0}.
\end{aligned}
\label{eq:general-problem}
\end{equation}

其中，所有场景都必须满足可行性要求，包括权重为零的场景。完成时间和温度作为约束，目标只优化完整系统能量。



一般问题保留连续启动时刻。为使热传播与能量积分所用的时间分段能够在求解前确定，下文采用预先声明的有限启动集合建立MILP参照。具体地，对每个实例和模式给定有限集合\(\mathcal G_{cim}\)，并定义

\begin{equation}
\mathcal F_{\mathcal G}=\{(u,m,s)\in\mathcal F:
s_{ci}\in\mathcal G_{ci,m_{ci}}\},\qquad
(P_{\mathcal G})\quad\min_{(u,m,s)\in\mathcal F_{\mathcal G}}\bar E.
\label{eq:finite-start-problem}
\end{equation}

\(P_{\mathcal G}\)只限制启动选择，使用与\(P\)相同的响应、功耗、温度和能量关系。集合可以由等距格点构成，也可以显式列出其他允许起点；时长和结束时刻均不取整，温度仍按连续时间约束。温度和能量等量在MILP中仍可用连续变量表示。



由于\(\mathcal F_{\mathcal G}\subseteq\mathcal F\)，限制起点后最小能耗只能保持或增加。有限域安全方案也是一般问题的可行方案，其最优值是一般问题最优值的上界；有限域无解不排除其他连续起点。下文围绕\(P_{\mathcal G}\)建立变量、约束和求解流程。



\section{问题分析与形式化}



共享部署把各场景联系在同一个优化问题中：某个Stage一旦部署到NPU0，其在所有场景中的实例都要使用NPU0，而频率和起点仍可分别调整。响应差异改变每种选择的占用时长，频率域忙闲变化及温度漏电又使能耗依赖完整时序。因此，形式化时需要同时保留部署、工作点、起点及热能耗关系。



为便于理解各项约束，表\ref{tab:constraint-example}给出一个场景内的示意安排。请求A包含依赖链\(A_1\rightarrow A_2\)，请求B只含\(B_1\)；二者均在0 ms释放，期限为6 ms，评价窗口为8 ms。CPU1与CPU2同属一个频率域，H表示该域的一个合法工作点，N表示NPU0的一个合法工作点。下列模式和时长仅用于说明约束，不作为实验输入，也不预设其能耗与热安全结果。



\begin{table}[htbp]
\centering
\caption{用于解释约束的场景内安排}
\label{tab:constraint-example}
\begin{tabular}{cccc}
\toprule
Stage & 执行单元 & 工作点 & 执行区间（ms） \\
\midrule
\(A_1\) & CPU1 & H & $[0,2)$ \\
\(A_2\) & NPU0 & N & $[2,5)$ \\
\(B_1\) & CPU2 & H & $[0,3)$ \\
\bottomrule
\end{tabular}
\end{table}



若启动时刻连续变化，功率片段的长度也随决策变化，热传播中的矩阵指数便依赖未知片长，不能直接作为普通MILP的常系数。有限启动集合提供了一种明确的处理方式：先由所有合法起止事件建立固定时间分段，再用二元变量选择实际发生的工作。这样，各片段长度已知，固定控制状态下的温度与能量传播可以写成仿射关系。



为此，对实例\(i\)建立选项集合

\begin{equation}
\begin{aligned}
\mathcal Q_{ci}=\{o=(i,m,t):\;&m\in\mathcal M_{ci},\ t\in\mathcal G_{cim},\\
&r_{c,r(c,i)}\le t,\quad t+d_{\tau(c,i)}(m)\le W_c\},\\
\mathcal Q_c&=\bigcup_i\mathcal Q_{ci}.
\end{aligned}
\label{eq:options}
\end{equation}

选项\(o\)具有已知单元\(u_o\)、域\(h_o=\gamma(u_o)\)、工作点\(\kappa_o\)、起点\(t_o\)、时长\(d_o\)、活动功率\(q_o\)及终点\(f_o=t_o+d_o\)。集合已删除违反释放或窗口必要条件的选项，其余声明选项均保留。这一步不排除任何有限域合法安排；建立的是单Stage选项，不预先枚举完整时间表。



将全部选项起点、真实终点以及\(0,W_c\)排序去重，得到

\begin{equation}
0=e_{c0}<e_{c1}<\cdots<e_{cL_c}=W_c,\qquad
\delta_{c\ell}=e_{c,\ell+1}-e_{c\ell}>0.
\label{eq:events}
\end{equation}

这些事件包含任何候选的活动或工作点改变时刻。对片段\([e_{c\ell},e_{c,\ell+1})\)，定义已知覆盖系数\(\chi_{co\ell}\)：选项覆盖该片段时为1，否则为0。选项终点即使不是启动格点，也作为独立事件保留。



固定事件集取自全部选项，而非仅取被选安排。表\ref{tab:constraint-example}中的被选项在0、2、3、5 ms发生起止事件；若\(A_1\)还允许在1 ms启动，则1 ms也进入固定事件集。即使最终仍选0 ms启动，1 ms仍保留为计算边界，却不触发新的工作、换点或温度重置。未选项增加分段，不增加实际执行。



\subsection{任务分配约束}



每个固定Stage只部署到一个执行单元。引入二元变量\(x_{\tau u}\)，取1表示固定Stage \(\tau\)部署到单元\(u\)。每个实例则通过二元变量\(y_{co}\)选择一个模式与起点选项。由于实例的单元必须与共享部署一致，有

\begin{equation}
\begin{aligned}
&\sum_{u\in\mathcal U}x_{\tau u}=1,
&&\tau\in\mathcal J,\\
&\sum_{o\in\mathcal Q_{ci}}y_{co}=1,
&&i\in\mathcal I_c,\ c\in\mathcal C,\\
&\sum_{\substack{o\in\mathcal Q_{ci}\\u_o=u}}y_{co}=x_{\tau(c,i),u},
&&i\in\mathcal I_c,\ u\in\mathcal U,\ c\in\mathcal C.
\end{aligned}
\label{eq:milp-assignment}
\end{equation}

没有合法选项的单元对应空和0，因此所有实例--单元链接均须保留。该表达共用映射，不要求不同场景的起点和工作点相同。



\subsection{截止时间约束}



选中选项后，Stage的启动和完成时刻分别为

\begin{equation}
s_{ci}=\sum_{o\in\mathcal Q_{ci}}t_oy_{co},\qquad
F_{ci}=\sum_{o\in\mathcal Q_{ci}}f_oy_{co}.
\label{eq:option-times}
\end{equation}

这些表达式均为线性式，执行时长已经通过模式响应计入\(f_o\)。



请求释放以前不能执行其Stage，后继Stage需等待前驱完成。因此

\begin{equation}
s_{ci}\ge r_{c,r(c,i)},\qquad
F_{ci}\le s_{cj}\quad ((i,j)\in\mathcal E_c).
\label{eq:precedence}
\end{equation}

每个请求的全部必需输出还需按时完成，全部实例应位于评价窗口内：

\begin{equation}
F_{ci}\le D_{cr}\quad(i\in\mathcal O_{cr}),\qquad
0\le s_{ci},\quad F_{ci}\le W_c.
\label{eq:deadlines}
\end{equation}

请求完成时间在输出时取\(\max_{i\in\mathcal O_{cr}}F_{ci}\)。逐输出约束可以表达该最大值不超过期限，无须为此增加一个自由完成时间变量。示例中，\(A_2\)在\(A_1\)结束后启动，请求A、B分别在5、3 ms完成，均满足6 ms期限；5--8 ms的空闲尾段仍进入后续热能耗评价。



\subsection{执行单元约束}



一个执行单元在任一片段内最多执行一个Stage。利用半开区间覆盖系数，可直接写为

\begin{equation}
\sum_{o\in\mathcal Q_c:\,u_o=u}\chi_{co\ell}y_{co}\le1,
\qquad u\in\mathcal U,\quad 0\le\ell<L_c.
\label{eq:occupancy}
\end{equation}

由于所有可能的起止时刻都在事件集合中，片段内活动集合不会变化，逐片段检查即覆盖完整资源占用。一个Stage结束与另一个启动重合时，二者不会同时覆盖同一片段。示例中\(A_1\)与\(B_1\)使用不同CPU单元，可以重叠；若同时分配到CPU1，重叠片段的占用和将为2，从而违反式\eqref{eq:occupancy}。不同单元上的并发还须满足下述频率域约束。



\subsection{工作点控制约束}



同域单元共享工作点，而域完全空闲时立即采用最低档$L_h$。引入二元变量\(\rho_{ch\ell k}\)，取1表示域\(h\)在片段\(\ell\)使用工作点\(k\)。每段只选一个状态：

\begin{equation}
\sum_{k\in\mathcal K_h}\rho_{ch\ell k}=1,
\qquad 0\le\ell<L_c.
\label{eq:domain-onehot}
\end{equation}

被选选项在全部执行片段中必须采用自身工作点，因此

\begin{equation}
y_{co}\le\rho_{c,h_o,\ell,\kappa_o},
\qquad \chi_{co\ell}=1.
\label{eq:active-lock}
\end{equation}

若两个同域选项重叠且工作点不同，式\eqref{eq:domain-onehot}与式\eqref{eq:active-lock}会禁止同时选择；同点并发则不受此限制。



为了避免控制域在没有活动Stage时停留在较高档位，对每个非最低工作点$k\ne L_h$增加活动支持条件：
\begin{equation}
\rho_{ch\ell k}\le
\sum_{\substack{o\in\mathcal Q_c:\,h_o=h,\ \kappa_o=k}}
\chi_{co\ell}y_{co},\qquad k\in\mathcal K_h\setminus\{L_h\}.
\label{eq:idle-auto-low}
\end{equation}
当控制域处于空闲片段时，式\eqref{eq:idle-auto-low}右侧为零，所有非最低档均不能被选；结合式\eqref{eq:domain-onehot}，该片段只能取$L_h$。当某个Stage正在执行时，式\eqref{eq:active-lock}要求该域采用Stage的工作点；若有其他同域Stage加入或退出，但原Stage尚未完成，工作点就不能改变。这样，忙段中途的事件、同刻结束与启动、首次启动前以及完成后的空闲尾段，都可由同一组条件描述，不需要旧版空闲保持上一档的约束。

温度系数可能同时依赖多个域的状态。记\(\Omega\)为输入给定的物理相容联合状态集合，\(\omega_h\)为其中域\(h\)的工作点，并用二元变量\(w_{c\ell\omega}\)关联各域选择：

\begin{equation}
\sum_{\omega\in\Omega}w_{c\ell\omega}=1,\qquad
\sum_{\omega\in\Omega:\,\omega_h=k}w_{c\ell\omega}=\rho_{ch\ell k}.
\label{eq:joint-state}
\end{equation}

\(\rho\)逐域说明工作点，\(w\)将同一片段各域的选择合成一个完整状态；式\eqref{eq:joint-state}保证两种表示一致。例如“CPU域为H、NPU域为N”是一个联合状态，不要求两个设备采用相同频率。联合状态决定各计量组的\(\xi_g(\omega)\)，输入须覆盖全部相容组合及入口组合。域空闲时，其在联合状态中的分量为最低档$L_h$；对应的基准功率和温度系数仍须进入功热计算。



\subsection{温度约束}



每段的入口温度由前一段的出口温度传递，当前工作点和活动集合再决定本段怎样升温或降温。固定联合状态\(\omega\)后，基准与漏电系数确定；记\(v_{c\ell}=(\chi_{co\ell}y_{co})_{o\in\mathcal Q_c}\)为片段内活动选择向量，则温度方程可整理为

\begin{equation}
\dot\theta=A_\omega\theta+b_\omega+D_cv_{c\ell},\qquad
A_\omega=C^{-1}\bigl[-G+\operatorname{diag}(\ell_\omega)\bigr],
\quad \ell_{\omega,n}=\sum_g\lambda_{gn}(\xi_g(\omega)).
\label{eq:state-ode}
\end{equation}

其中，\(b_\omega\)包含分配后的基准和参考温度偏置，\(D_c\)将选项活动功率注入对应节点，具体定义见章末补充推导。对已知片长\(\delta_{c\ell}\)积分，得到

\begin{equation}
w_{c\ell\omega}=1\ \Longrightarrow\
\theta_{c,\ell+1}=\Phi_{c\ell\omega}\theta_{c\ell}
+\psi_{c\ell\omega}+\Gamma_{c\ell\omega}v_{c\ell}.
\label{eq:thermal-transition}
\end{equation}

这里\(\theta_{c\ell}\)与\(\theta_{c,\ell+1}\)分别是本段入口和出口温升，\(\Phi\)传播已有热状态，\(\psi\)表示基准与参考偏置的影响，\(\Gamma v\)表示活动工作的影响。矩阵指数及积分只涉及给定系数和固定片长，所以右端对变量是仿射的，不含两个连续决策变量的乘积。



条件等式可用有界线性约束实现。例如，令温度等式残差为\(R^T_{c\ell\omega,n}\)，则写成

\begin{equation}
-M^T_{c\ell\omega,n}(1-w_{c\ell\omega})
\le R^T_{c\ell\omega,n}
\le M^T_{c\ell\omega,n}(1-w_{c\ell\omega}).
\label{eq:thermal-bigM}
\end{equation}

当\(w=1\)时，上下界都为零，残差必须为零，所选状态的传播等式生效；当\(w=0\)时，该状态的等式不应限制实际轨迹。因此，\(M^T\)必须覆盖未激活行的全部可能残差。它是条件表达的松弛界，不是允许超温的容差。章末根据有限窗口和已知系数给出保守有限界。



入口温度与事件端点温限为

\begin{equation}
\theta_{c0}=T_c^0-T_a\mathbf1,\qquad
T_a+\theta_{c\ell n}\le T_n^{\max},\quad 0\le\ell\le L_c.
\label{eq:event-temperature}
\end{equation}

窗口末点$W_c$没有执行中的Stage，应按各域最低档的入口状态核查该端点的功率有效性；这不增加额外的积分区间。

同时加入声明有效温区及物理有效条件的端点必要约束。端点合格还不足以保证事件之间不超温，因此这些条件构成连续热要求的放松。完整片段检查与后续温限行在下一节连接。



\subsection{目标函数与优化问题}



完整能耗需要计入温度变化引起的功率变化。固定状态\(\omega\)时，边界内电功率与温度、活动选择的关系为

\begin{equation}
P^{\rm net}(t)=\beta_\omega+\ell_\omega^{\mathsf T}\theta(t)
+q_c^{\mathsf T}v_{c\ell},\qquad q_c=(q_o)_o.
\label{eq:state-power}
\end{equation}

将式\eqref{eq:state-ode}的温度解代入并积分，可得片段电能的仿射表达

\begin{equation}
w_{c\ell\omega}=1\ \Longrightarrow\
\Delta E_{c\ell}=\mu_{c\ell\omega}^{\mathsf T}\theta_{c\ell}
+\zeta_{c\ell\omega}+\nu_{c\ell\omega}^{\mathsf T}v_{c\ell}.
\label{eq:energy-transition}
\end{equation}

这条条件等式同样采用有效残差界线性化。\(\Delta E\)已经包含基准、温度修正和活动功率，不再额外添加一份\(q_od_o\)。无Stage活动的片段仍传播温度并计能，因而等待和尾段均进入目标。



记给定边界外能量为\(E_{{\rm ext},c}=\int_0^{W_c}P_{{\rm ext},c}(t)dt\)，当前已加入的连续热必要行集合为\(\mathcal H^{(j)}\)。第\(j\)轮主问题为

\begin{equation}
\begin{aligned}
\mathcal M^{(j)}:\quad
&\min\ \sum_{c\in\mathcal C}p_c
\left(E_{{\rm ext},c}+\sum_{\ell=0}^{L_c-1}\Delta E_{c\ell}\right)\\
&\text{s.t. 式\eqref{eq:milp-assignment}--\eqref{eq:joint-state}的调度与控制约束，}\\
&\hspace{1.5em}\text{式\eqref{eq:thermal-transition}、\eqref{eq:energy-transition}的条件线性形式，}\\
&\hspace{1.5em}\text{入口、端点、有效域必要约束及有依据的变量界，}\\
&\hspace{1.5em}\mathcal H^{(j)},\quad x,y,\rho,w\in\{0,1\}.
\end{aligned}
\label{eq:master-problem}
\end{equation}

温度、段能量及起止表达使用连续变量；其余系数在输入给定后计算。完整能耗直接参与部署、工作点和时序选择，外部评价器复核同一个目标，而不是用活动能量规划后再另行排名。



所有声明有限选项都能进入该表达。反过来，任一满足整数约束的选项组合也确定一份资源、期限及控制合法的时间表；给定入口后，热能耗轨迹唯一。对应关系和条件等式的系数推导见章末。\(\mathcal M^{(j)}\)仍可能含有片段内部超温的安排，下面通过连续检查逐步补充必要条件。



上述固定片长下的矩阵指数、积分与条件线性化描述精确实数意义下的模型关系。实际生成MILP系数时需考虑数值传播误差，工程实现对可核验误差采用向外放松与有效包络，因此求解器内部温度、能量和界并不能未经检查就替代真实电热轨迹。完整安全性及最终能量仍由独立连续评价器按原模型验证，数值主问题的报告状态与可信安全能量上界分别保留。

\section{MILP与连续热反馈求解流程}



\subsection{方法概述}



上一节的MILP已经包含温度传播、事件端点温限和完整能耗目标，可据此共同选择部署、工作点和时序。尚需补齐的是事件之间的连续安全检查。为此，独立检查器按候选恢复完整执行过程，复核温度、物理有效范围和能量；安全候选形成可用方案，明确违规则形成反馈约束，使规划器重新选择配置或时序。



该组织将两类计算分开：固定事件下的调度、状态递推和能量积分进入线性主问题；任意连续时刻的温度极值由外部检查器处理。两者使用相同输入、功耗模型和热方程。图\ref{fig:solver-flow}给出主要信息流。



\begin{figure}[htbp]
\centering
\begin{tikzpicture}[
  box/.style={draw,rounded corners,align=center,text width=7cm,inner sep=6pt},
  >=Stealth,node distance=10mm,font=\small]
\node[box] (input) {输入场景、平台、功耗与热参数\\建立有限选项、固定事件和状态系数};
\node[box,below=of input] (milp) {求解或续解完整能耗MILP\\共同选择共享部署与各场景执行安排};
\node[box,below=of milp] (check) {独立复核全部场景\\资源、期限与完整能量\\连续温度与物理有效范围};
\node[box,below=of check] (record) {汇总已有安全方案、界与状态\\明确违规时形成有效约束\\判断继续或停止};
\node[box,below=of record] (output) {输出已有方案、能量与证据范围\\区分已证结论和未决事项};
\draw[->] (input)--(milp);
\draw[->] (milp)--node[right]{取得候选}(check);
\draw[->] (check)--(record);
\draw[->] (record)--node[right]{已证结论或未决停止}(output);
\draw[->] (record.west)--++(-1.1,0)|-
node[pos=0.38,left,align=right]{明确违规后反馈\\或安全未闭合续解\\且允许继续}(milp.west);
\draw[->] (milp.east)--++(1.1,0)|-
node[pos=0.38,right,align=left]{无解或停止\\且无待查候选}(record.east);
\end{tikzpicture}
\caption{完整能耗MILP与连续热检查的信息交互}
\label{fig:solver-flow}
\end{figure}



\subsection{连续热反馈的引入}



连续温限要求每个时刻都不超过上限，而不是只要求任务起止时刻安全。即使一个片段内工作点及活动状态不变，各节点也可能因导热和漏电反馈先升温后降温，形成内部峰值。有限端点约束不能预先覆盖所有这类位置。



对一个完整候选，检查器先由选项恢复时间表和工作点与空闲最低档轨迹，再逐片段评价式\eqref{eq:state-ode}。每个节点的温度在片段内部可微，峰值由端点或内部驻点取得；实际计算须用能覆盖整个区间的极值检查或保守误差包络，不能用密集采样代替连续结论。各场景从各自入口检查到\(W_c\)，包括没有活动的等待与尾段。



若场景\(c\)在片段\(\ell\)、局部时刻\(\tau^*\in[0,\delta_{c\ell}]\)的节点\(n\)明确超温，则在同一完整控制状态\(\omega\)下，该时刻的温度可写为

\begin{equation}
\widehat\theta(\tau^*)=
e^{A_\omega\tau^*}\theta_{c\ell}
+\mathsf K_\omega(\tau^*)\bigl(b_\omega+D_cv_{c\ell}\bigr),
\quad \mathsf K_\omega(t)=\int_0^t e^{A_\omega z}\,dz.
\label{eq:interior-temperature}
\end{equation}

由全部安全安排在该时刻都必须满足温限，可加入

\begin{equation}
w_{c\ell\omega}=1\ \Longrightarrow\
\mathbf e_n^{\mathsf T}\widehat\theta(\tau^*)
\le T_n^{\max}-T_a.
\label{eq:conditional-cut}
\end{equation}

\(\mathbf e_n\)为取节点\(n\)分量的单位向量。反馈固定了片段、检查时刻与适用状态，未固定产生该温度的整张时间表：\(\theta_{c\ell}\)仍由之前的安排决定，\(v_{c\ell}\)仍可变化。因此，其他具有相同当前控制状态的候选也须满足该行；若候选改用其他联合状态，该行不激活，其热轨迹按对应系数重新检查。



例如，若某候选在一个CPU与NPU共同工作的片段内部发生NPU超温，反馈不会指定“把哪个Stage延迟多久”。规划器仍可减少之前的热积累、改变当前重叠，或选择其他工作点和部署。反馈提供必须满足的温限，采用哪种调整由完整能耗目标与其余约束共同决定。



当明确违规已经确认、但无法可靠形成共享温限行时，可使用完整赋值排除。设\(S_c^*\)为该场景选中的全部选项，则

\begin{equation}
\sum_{o\in S_c^*}y_{co}\le |\mathcal I_c|-1.
\label{eq:full-assignment-cut}
\end{equation}

该式只删除当前完整模式与起点组合，同一配置的其他时序仍被保留。由于场景入口独立，其他场景的安排不能改变此完整赋值的违规事实。物理有效域的明确违规可按相同原则反馈；温度或有效性尚未分辨时，不添加排除。



\subsection{具体求解流程}



\textbf{（1）初始化与主问题构造。} 检查输入、参数有效范围及注热守恒，建立选项、固定事件和相容联合状态。计算温度传播、能量积分系数及有效变量界，加入资源、期限、控制、端点温限等约束，得到初始\(\mathcal M^{(0)}\)。这一步不读取其他方法的候选方案或参考答案。



\textbf{（2）主问题求解。} MILP以完整加权能耗为目标搜索候选。整数解解码后，先按原始服务条件复核共享部署、依赖、资源互斥以及共享频率域与空闲最低档规则。若候选不符合这些条件，应作为数值或实现异常处理，不能据其生成热排除。



\textbf{（3）连续评价。} 对全部场景独立重建温度和能量，检查完整窗口中的温限与模型有效范围。只有所有场景均完整、合法且安全时，才汇总加权能量并更新当前安全上界\(U\)。零权重场景仍检查安全，不能用某场景未完成的前缀能量组成有效结果。安全但尚未证最优的候选保留为当前可行解，由同一主问题继续优化；若求解接口不能继续，则保留结果并报告最优性未决。



\textbf{（4）反馈与继续。} 对明确违规优先形成式\eqref{eq:conditional-cut}的有效必要行，必要时使用式\eqref{eq:full-assignment-cut}；随后在剩余预算内重新求解。新约束约束完整能耗目标下的候选选择，因此反馈既可引起起点调整，也可引起工作点或共享部署重选。一次时序失败不直接排除整个配置。



\textbf{（5）停止与输出。} 若主问题的全局最优候选通过连续检查，则在精确系数和有效反馈条件下，它也是\(P_{\mathcal G}\)的最优解；若保留原有限域全部安全方案的主问题被可靠证明无解，则该有限域无可行解。只有安全候选但界未闭合时，输出安全方案及未决最优性。数值未分辨、预算结束或资源限制时，保留已有安全方案、可靠界及未决原因，不把停止本身解释为无解。



表\ref{tab:solver-outline}汇总上述流程。所有建模、系数计算、重复求解与检查成本均应计入方法预算，不能每轮重新获得完整预算。



\begin{table}[htbp]
\centering
\caption{MILP与连续热反馈的求解步骤}
\label{tab:solver-outline}
\begin{tabular}{p{0.06\textwidth}p{0.85\textwidth}}
\toprule
步骤 & 处理及信息保留 \\
\midrule
1 & 读入模型与有限启动域，建立初始主问题，初始化安全记录与反馈集合。\\
2 & 在剩余预算内求解，保存内部目标、界、候选和退出状态。\\
3 & 解码候选；复核调度合法性，并对全部场景作完整热能耗评价。\\
4 & 全部安全：更新\(U\)；可信闭合则结束，否则继续当前主问题的优化，或以最优性未决结束。\\
5 & 明确违规：加入有效必要行或完整赋值排除，返回步骤2。\\
6 & 未分辨或无法继续：保存已有证据，以对应未决状态结束。\\
\bottomrule
\end{tabular}
\end{table}



上述最优性结论来自主问题与连续问题的集合关系，而不是求解器退出标志本身。有限热行保留全部安全安排，因此精确求解的主问题最优值是有限域原问题的下界；其最优候选一旦安全，就同时达到该下界。若每次明确违规都排除了当前完整赋值，有限候选域在精确求解和可判定检查的理想条件下可在有限轮结束。



实际实现还需分别控制传播系数、积分系数、线性约束及连续检查的数值误差。若安全方案的场景能量评价为区间\([E_c^-,E_c^+]\)，应以\(\sum_cp_cE_c^+\)更新上界\(U\)，而非使用未经误差核查的名义值。候选经独立检查安全，并不自动证明浮点主问题没有误删可行方案，或其内部界是独立数学证书。因此，安全方案及其能量范围、可信上下界、求解器报告的界与状态分列保存；本章的正式实验采用统一的连续安全判据，分别保存实际安全上界、可信界与求解器内部状态；下一节据此报告已观察到的结果。



本节描述的主问题、热检查及有效反馈是精确模型下允许的求解关系。正式工程实现以同一次方法的实际剩余时间和返回状态决定能否继续求解；如果已经取得完整安全方案却没有可继续处理的MILP整数候选，也会保留该安全方案并正常报告最优性未决。不能把上述数学迭代关系直接理解为每次正式调用都执行了温限反馈。在P0、P1、P2正式实验中，外层MILP没有形成新增温限反馈行，下一节将分别展示安全检查、候选接纳和预算内搜索质量的实际记录。

\section{仿真结果与分析}
\label{sec:experiments}

本章的联合规划模型同时刻画了计算阶段的执行响应、共享频率域中的实际功率分布，以及连续变化的芯片温度。为分析这些因素对调度决策和系统能耗的影响，本节使用预先固定的抽象异构平台与请求实例开展仿真。实验并不只比较最终能量大小，而是分别回答三个问题：任务的频率响应信息能否改变调度选择；在执行模式不变的情况下，完整系统能量为何还会随时序和温度变化；连续温限是否会影响候选方案的接纳。在此基础上，再考察有限规划预算下不同方法实际得到的安全方案质量与计算开销。

为形成统一的比较基础，实验设置了整体MILP、两阶段固定设备联动局部MILP、HEFT适配版和模拟退火（Simulated Annealing，SA）四类方法。经典HEFT首先根据向上优先级确定DAG节点的安排顺序，再以最早完成时刻选择处理器\cite{topcuoglu2002heft}。本实验采用的HEFT适配版在此基础上处理跨场景的共同Stage身份、具体设备选择和各场景工作点，并进行一遍以完整安全能量严格下降为接纳条件的设备、工作点及起点调整。因此，后文的HEFT数据不是未经修改的经典HEFT算法结果。SA采用模拟退火的接受与扰动思想\cite{kirkpatrick1983sa}，从合法配置出发搜索完整安排，其每次候选仍按本章的资源、期限和电热安全条件进行评价。

整体MILP对应第3.5节范围较广的联合决策模型；两阶段联动方法则从本次调用自身构造的一份完整安全安排出发，固定各Stage的具体执行设备，依次调整不同Stage集合的工作点和起点。后者第二阶段只在前一阶段留下的固定配置条件下开放另一部分决策，属于原联合问题的受限搜索范围。各方法使用相同的任务与物理输入和安全评价标准，但搜索空间与构造机制并不相同。本节将其作为有限预算内的结果对照，第四章再详细分析这些算法的内部组织。

\subsection{仿真参数设置与实验设计}

仿真平台包含CPU1、CPU2、GPU0和NPU0四个执行单元，其中CPU1与CPU2共享一个CPU频率控制域，GPU0与NPU0各自具有独立的控制域。三个控制域均提供L、M、H三个离散电压--频率工作点，并在完全空闲时立即回到L档。不同处理器对Stage的合法支持、运行时长与活动功率分别设定，例如A的主计算阶段能够在CPU、GPU或NPU上执行，B的主计算阶段不支持GPU，C的主计算阶段不支持NPU。CPU域的基准功率只计一次，两个CPU单元的温度分别由各自的热节点记录。

表\ref{tab:exp-platform}给出主要频率及最低档基准，完整的三档电压、基准功率、温度系数、九个Stage的响应参数和四节点RC矩阵列于工程的实验输入对照文件，并可追溯至冻结配置。全部功热参数均为文献指导下声明的抽象仿真输入，而不是针对某一真实移动芯片联合测量所得的数据。环境温度与漏电参考温度为25\,$^\circ$C，模型有效温区为20--100\,$^\circ$C；系统还有0.30\,W的公共电功率，计入完整能耗但不注入当前四节点热网络。

\begin{table}[htbp]
\centering
\caption{端侧异构处理器的主要仿真工作点}
\label{tab:exp-platform}
\begin{tabular}{lccc}
\toprule
频率控制域 & L/M/H频率（GHz） & 最低档基准功率（W） & 执行单元 \\
\midrule
CPU & 0.8/1.2/1.6 & 0.60 & CPU1、CPU2 \\
GPU & 0.4/0.6/0.8 & 0.40 & GPU0 \\
NPU & 0.5/0.75/1.0 & 0.25 & NPU0 \\
\bottomrule
\end{tabular}
\end{table}

应用集合为A、B、C。在此次实验中，每个应用都由预处理、主计算与后处理三个串行Stage构成。系统分别考虑A+B和A+C两个场景，每个场景的每个应用各有一个请求，因而每个场景包含两条三Stage链、六个执行实例。应用A的固定Stage在两个场景中必须部署到同一具体执行单元；两个场景根据各自的请求和竞争条件分别确定工作点与启动时刻。各请求的释放时刻为0\,ms，场景权重均取$1/2$。

为了在不改变图结构的情况下形成不同的计算负载，实验对三个应用分别采用参考执行时长缩放系数。原H参考模式保持Stage及设备相关的响应参数，R1--R5仅改变A、B、C的参考时长比例，具体如表\ref{tab:exp-workload}所示。三Stage链的请求截止期限则根据该应用\emph{缩放后的CPU最高档参考执行时间总和}构造：
\begin{equation}
 D_{cr}=r_{cr}+\beta\,C_{a(r),{\rm CPU}}^{H},
 \qquad \beta=1.3,
 \label{eq:exp-deadline-generation}
\end{equation}
其中$C_{a(r),{\rm CPU}}^{H}$表示该应用三个Stage在CPU最高工作点上的缩放后参考时长之和。期限不由被评价算法的输出反向生成。因此，R1--R5是五个\emph{工作量比例变体}，并不是五种请求数量、DAG拓扑或不同期限系数的实验。

\begin{table}[htbp]
\centering
\caption{五组冻结工作量变体的应用时长缩放因子}
\label{tab:exp-workload}
\begin{tabular}{lccc}
\toprule
变体 & 应用A & 应用B & 应用C\\
\midrule
R1 & 1.0 & 1.0 & 1.0\\
R2 & 0.8 & 1.2 & 1.0\\
R3 & 1.2 & 0.8 & 1.0\\
R4 & 1.0 & 0.8 & 1.2\\
R5 & 1.0 & 1.2 & 0.8\\
\bottomrule
\end{tabular}
\end{table}

每个场景的能量评价窗口固定为400\,ms，合法Stage启动时刻取5\,ms网格，实际持续时间及结束时刻不作网格取整。Stage运行期间不抢占、不迁移，额外的数据搬运、交接与频点切换时延和能量在输入中设为零。功率计量同时覆盖Stage活动、等待及运行结束后的空闲尾段；连续热安全检查则覆盖完整窗口内所有节点的温度极值。

为区分执行响应与入口热条件，正式实验设置了表\ref{tab:exp-groups}中的P0、P1、P2三组输入。P0采用Stage和设备相关的原H响应、50\,$^\circ$C初温；P1保持其他参数，仅将所有方法共同使用的\emph{实际执行响应}改为统一反比模型R（$\alpha=1$）；P2恢复H响应，仅将场景入口温度提高到78\,$^\circ$C。各组均保持80\,$^\circ$C的温限与相同的合法控制策略。因此，P0与P1的能量数字属于不同真实执行时间表，不能直接用于估计响应信息准确性带来的节能；比较规划信息本身的价值，还需要额外的同真实H对照。

\begin{table}[htbp]
\centering
\caption{正式仿真的三组输入条件}
\label{tab:exp-groups}
\begin{tabular}{lcccc}
\toprule
组别 & 真实执行响应 & 初始温度 & 温度上限 & 输入变体 \\
\midrule
P0 & 阶段相关响应H & 50\,$^\circ$C & 80\,$^\circ$C & R1--R5 \\
P1 & 统一反比响应R & 50\,$^\circ$C & 80\,$^\circ$C & R1--R5 \\
P2 & 阶段相关响应H & 78\,$^\circ$C & 80\,$^\circ$C & R1--R5 \\
\bottomrule
\end{tabular}
\end{table}

本节按以下四个实验问题组织分析。实验一固定真实H响应，在R1、R2两个输入上比较使用H和统一R信息所作的规划，两类方案均在同一真实H模型中接受后验评价；它用于检验频率响应信息对模式与时序选择的影响。实验二使用相同Stage模式、仅改变合法启动时刻的数学见证，检验共享频率域与完整窗口能量之间的联系。实验三比较同一排程在P0与P2不同入口温度下的能量组成，并结合温度违例和独立压力条件分析连续热安全的接纳作用。实验四汇总P0、P1、P2中四类求解方法的安全能量和耗时，为下一章的高效算法设计提供实际问题依据。

实验四中，每个输入运行整体MILP、两阶段固定设备联动局部MILP、HEFT适配版各一次，以及五个独立随机种子的SA调用，故每组共有$5\times(1+1+1+5)=40$次方法调用，三组合计120次。各次方法的总计算预算均为300\,s，包括输入处理、初始化、MATLAB启动、构模、求解以及内部候选安全评价。每份最终安排还接受一次独立完整窗口安全复核，事后复核不继续搜索或替换预算内能量。实验一的H02受控对照不属于这120次：其中每次规划预算为120\,s，最终候选的真实H独立复核上限为60\,s，四次规划均使用相同的HEFT适配机制。

全部实验主要观察三项指标：按两场景权重得到的完整系统能量$U$（J），按期限、资源及连续温度要求复核后的安全性，以及从方法输入处理起算的实际方法根时钟（s）。对于已经完成独立评价的安全安排，以可信能量区间上端作为其预算内安全能量$U$；比较两种方法时，正的
\begin{equation}
 R_E=\frac{U_{\rm HEFT}-U_{\rm method}}{U_{\rm HEFT}}\times100\%
 \label{eq:exp-saving}
\end{equation}
表示被比较方法在这一输入中能量更低。严格大小由原归档中的高精度区间判断，表中的有限小数仅便于阅读。SA的五个随机种子是五次独立方法调用，其能量中位数只作为描述性汇总，并不代表一个算法在单次300\,s内进行了五次完整搜索。300\,s是部署前的计算预算，400\,ms是被规划推理服务的评价窗口；表内能量不计入规划器自身运行所消耗的电能。


\subsection{任务响应与共享部署的影响}

实验一首先考察频率响应信息是否会改变调度选择。根据式\eqref{eq:response-model}，相同频率下降比例对不同Stage的持续时长产生不同影响；在共享CPU频率域和跨场景共同设备部署的条件下，某个Stage延长执行还可能影响同时运行的其他Stage，进而改变可行的工作点和时序。

P0和P1采用不同的实际响应模型，二者只能用于观察真实输入变化后的调度表现。为了进一步把\emph{规划信息变化}与\emph{实际物理输入变化}区分开来，H02对照在R1和R2上始终保留原H响应作为真实执行参数：一组HEFT适配版规划直接使用H响应信息，另一组仅在规划阶段使用统一反比R信息，最后四份方案都根据同一真实H的功热、期限及安全规则进行独立复核。这种安排使得同一输入上的两份能量可以直接比较，而不需要将不同真实任务时长混在一起。

表\ref{tab:exp-response}列出了完整窗口加权能量。R1的H信息规划方案为0.8810879973\,J，统一R信息规划方案为0.9173594769\,J，前者相对后者低3.953900\%。两份方案的具体执行单元映射相同，但A+B场景的B预处理在H方案中选择CPU1的L档，而R方案选择H档；B主计算分别选择NPU0的M档与H档，其他Stage的起点也随之调整。已保存的分项差异表明，R方案的活动能量、域基准能量和温度修正能量均高于H方案。因此，任务相关响应信息在R1中改变了工作点和时序，并最终影响完整系统能耗。

\begin{table}[htbp]
\centering
\caption{相同真实H模型下不同规划响应信息的完整系统能量}
\label{tab:exp-response}
\begin{tabular}{lccc}
\toprule
输入 & H信息规划（J） & R信息规划（J） & H相对R的能量变化 \\
\midrule
R1 & 0.8810879973 & 0.9173594769 & 降低3.953900\% \\
R2 & 0.9554599350 & 0.9554455369 & 增加0.001507\% \\
\bottomrule
\end{tabular}
\end{table}

R2得到的结论不同：H信息与R信息方案的设备映射及工作点全部相同，仅A+C场景中四个Stage的启动时刻存在差异。R信息方案的真实H能量低约14.398\,$\mu$J，约占R信息方案能量的0.001507\%；该微小差异来自温度修正积分。这里没有出现H信息逐题节能的规律，说明在有限搜索预算下，更符合实际时长的输入可能改变搜索路径，却不保证该启发式在每个问题上都得到最低的真实能量。

两个输入还具有相同的跨场景部署规则：应用A的固定Stage在A+B和A+C中必须使用同一具体执行单元，但可以分别选择工作点和起点。R1的设备映射并没有随规划信息改变，所以其3.953900\%的能量差不能归因于跨场景映射变更。现有实验也没有在相同物理条件下移除共享部署要求，故这些结果支持的是\emph{既定共享部署约束下频率响应信息的条件性决策价值}，并不能证明共享部署本身比逐场景独立部署更省电。

\subsection{热能耗评价与连续反馈的作用}

实验二考察一个更基本的问题：即使所有Stage都使用相同的设备和工作点，完整系统能量是否仍取决于启动时刻。按式\eqref{eq:full-energy}，系统能量不仅包括Stage活动功率，还包括不同工作点的域基准以及温度相关漏电。对共享CPU域，两个高档Stage并行运行和先后接力虽然各自占用相同的高档计算时长，却可能产生不同的\emph{高档驻留时间并集}。

以P0/R1的冻结模式参数为例，构造两张完整合法、连续热安全的时间表P与S。两张表在A+B、A+C两种场景中选用完全相同的Stage模式与共享设备部署，A+C中的启动时刻也相同；仅将A+B中B预处理的起点从0\,ms调整为10\,ms。A与B的预处理分别在CPU1、CPU2的H档运行10\,ms和15\,ms。P中两者并行，共享CPU域H档占用区间为$[0,15)$\,ms；S中两者接力，H档区间为$[0,25)$\,ms。受空闲时即时回最低档L的规则影响，S比P额外支付10\,ms的高档域基准差。

由于CPU域H档和L档基准分别为1.30\,W和0.60\,W，A+B场景的参考基准能量差为
\begin{equation}
 E^{AB}_{{\rm base},S}-E^{AB}_{{\rm base},P}
 =(1.30-0.60)\times0.010=0.007\ {\rm J}.
 \label{eq:exp-C19-baseline}
\end{equation}
A+C场景保持不变，故两场景等权时直接基准差为3.5\,mJ。进一步在原四节点RC模型上考虑后续的温度漏电变化，可以得到完整窗口加权能量的严格解析界
\begin{equation}
 \bar E_S-\bar E_P>0.003254\ {\rm J}.
 \label{eq:exp-C19-complete}
\end{equation}
这两张构造时间表均满足请求期限和连续热安全条件。式\eqref{eq:exp-C19-complete}说明，同一组Stage模式向量并不足以唯一确定完整系统能耗，实际时间表和共享频率域的忙闲轨迹仍须进入目标函数。这里的数值来自原参数下的静态数学见证，不是C方法、HEFT或其他算法真实运行时额外节省了3.254\,mJ的结果。

实验三将注意力转向入口温度。在P0与P2中，R1的最终HEFT适配版时间表具有完全相同的具体设备、工作点和启动时刻；两组唯一改变的热条件为入口温度由50\,$^\circ$C提高至78\,$^\circ$C。在相同时间表下，基准、活动与公共能量保持不变，温度修正能量却由0.101350965\,J增加到0.207116548\,J，使完整系统能量增加约0.105765583\,J。表\ref{tab:exp-thermal-energy}给出了具体分项，图\ref{fig:exp-thermal-parts}进一步展示温度修正部分在总能量中的变化。

\begin{table}[htbp]
\centering
\caption{同一R1时间表在不同初始温度下的能量分项}
\label{tab:exp-thermal-energy}
\begin{tabular}{lcc}
\toprule
能量组成 & P0（50\,$^\circ$C，J） & P2（78\,$^\circ$C，J） \\
\midrule
共享域与设备基准 & 0.503150000 & 0.503150000 \\
Stage活动能量 & 0.156587033 & 0.156587033 \\
公共电能 & 0.120000000 & 0.120000000 \\
温度修正能量 & 0.101350965 & 0.207116548 \\
\midrule
完整系统能量 & 0.881087997 & 0.986853581 \\
\bottomrule
\end{tabular}
\end{table}

\begin{figure}[htbp]
\centering
\begin{tikzpicture}[x=1cm,y=5.6cm,font=\small]
 \draw[->] (0.85,0)--(0.85,1.08) node[above,font=\scriptsize] {能量/J};
 \foreach \y/\lab in {0/0,0.2/0.2,0.4/0.4,0.6/0.6,0.8/0.8,1.0/1.0}{
   \draw (0.79,\y)--(.9,\y);
   \node[anchor=east,font=\scriptsize] at (.72,\y) {\lab};
 }
 \fill[gray!46] (1.7,0) rectangle (2.7,.503150);
 \fill[blue!25] (1.7,.503150) rectangle (2.7,.659737033);
 \fill[orange!25] (1.7,.659737033) rectangle (2.7,.779737033);
 \fill[red!35] (1.7,.779737033) rectangle (2.7,.881087998);
 \fill[gray!46] (3.7,0) rectangle (4.7,.503150);
 \fill[blue!25] (3.7,.503150) rectangle (4.7,.659737033);
 \fill[orange!25] (3.7,.659737033) rectangle (4.7,.779737033);
 \fill[red!35] (3.7,.779737033) rectangle (4.7,.986853581);
 \draw[thin] (1.7,0) rectangle (2.7,.881087998);
 \draw[thin] (3.7,0) rectangle (4.7,.986853581);
 \node[above,font=\scriptsize] at (2.2,.891) {0.881088};
 \node[above,font=\scriptsize] at (4.2,.995) {0.986854};
 \node[below] at (2.2,-.03) {P0：50$^\circ$C};
 \node[below] at (4.2,-.03) {P2：78$^\circ$C};
 \fill[gray!46] (5.9,.88) rectangle (6.17,.92);
 \node[anchor=west,font=\scriptsize] at (6.36,.90) {域与设备基准};
 \fill[blue!25] (5.9,.72) rectangle (6.17,.76);
 \node[anchor=west,font=\scriptsize] at (6.36,.74) {Stage活动能量};
 \fill[orange!25] (5.9,.56) rectangle (6.17,.60);
 \node[anchor=west,font=\scriptsize] at (6.36,.58) {公共电能};
 \fill[red!35] (5.9,.40) rectangle (6.17,.44);
 \node[anchor=west,font=\scriptsize] at (6.36,.42) {温度相关能量};
 \node[anchor=west,font=\scriptsize] at (5.9,.17) {同一设备、工作点与时间表};
\end{tikzpicture}
\caption{同一R1时间表在两种入口温度下的完整能量组成}
\label{fig:exp-thermal-parts}
\end{figure}

图中两个柱形的基准、活动及公共部分高度完全相同，增加的部分来自温度修正。这一对照说明，若忽略入口热状态与温度漏电，系统能量评价就会丢失相同执行安排在不同热入口下的代价差异；它仍属于声明电热模型下的结果，不是实测芯片功耗或温度感知算法的净节能量。

温度限制还会改变候选方案能否被接纳。P0与P2曾保存五组完全相同的HEFT中间候选：这些候选在50\,$^\circ$C入口下都满足连续安全条件，而在78\,$^\circ$C入口下均出现明确超80\,$^\circ$C的节点温度。以R1为例，同一候选在P0的A+B场景CPU1峰温约为55.558515\,$^\circ$C，在P2则于95\,ms附近达到约81.209166\,$^\circ$C。由于后者超过既定上限，检查器拒绝了这一候选。P2的40次方法调用中，有17次保存了明确的热违规，共对应46条在各次调用内部去重的违规评价记录；但最终40份安排均通过完整连续温度安全复核，节点峰温上界最大约为79.858549\,$^\circ$C。

另一项独立压力实验将入口保持在50\,$^\circ$C，把温限从80\,$^\circ$C收紧为52\,$^\circ$C。该条件下原H档初始化在A+B的NPU和A+C的GPU节点出现明确超温，构造器转而尝试M档并获得完整安全安排。这说明温限要求会影响某些初始候选的模式选择与安全接纳。不过，P0与P2五组最终HEFT时间表仍然相同，不能将中间候选的热拒绝表述为最终部署已因初温变化而改变。

最后需要区分连续安全\emph{检查}与MILP外层\emph{反馈}。前者对每份候选时间表重建完整热轨迹，判断其是否安全；后者是在MILP候选被明确判定违规后，生成新的必要约束并继续求解。P2中的中间候选热拒绝属于HEFT和SA的内部评价，52\,$^\circ$C压力实验属于初始化模板的安全选择；三组共120次正式方法调用中的MILP外层连续热反馈新增行数均为零。因此，已有结果证明了热状态影响能耗与安全准入，但没有形成“反馈约束改变MILP搜索并产生额外节能或加速”的实验记录。


\subsection{求解表现与适用范围}

实验四对比各方法在共同模型、统一输入和指定方法预算下的最终安全能量与实际计算开销。三组正式实验各包含R1--R5五种工作量变体，每种变体独立执行整体MILP、两阶段固定设备联动局部MILP、HEFT适配版和五次SA搜索。全部120次方法调用均保留了一份满足资源、期限和完整连续热约束的安全安排，并完成规定的一次独立后验检查；这里的“安全”仅表示在声明的模型与评价窗口内可用，不表示已经获得某一完整问题的最优解。

在解释表\ref{tab:exp-energy}前，需要明确其中“整体MILP”列的来源。对P0、P1、P2的15个对应输入，整体MILP方法都实际启动了主问题构造与求解器，但在各自300\,s预算内\emph{没有收到可以独立完整评价的整数候选}；这15份最终安全安排均由该方法自身的初始化构造器形成，随后被保留并通过连续安全复核。因此，该列是“整体MILP方法调用在规定预算内保留的初始安全能量”，不是MILP整数搜索获得的改进解，更不是完整联合问题的数学最优值。它与联动方法的差异反映实际程序在有限时间内得到的方案质量，不意味着受限搜索域在数学上优于完整搜索域。

\begin{table}[htbp]
\centering
\small
\caption{P0、P1、P2三组实验的预算内完整安全能量（J）}
\label{tab:exp-energy}
\begin{tabular}{llrrrr}
\toprule
组别 & 变体 & 整体MILP & 联动MILP & HEFT & SA中位数 \\
\midrule
P0 & R1 & 1.084886 & 0.860663 & 0.881088 & 0.904844 \\
   & R2 & 1.078187 & 0.859268 & 0.955460 & 0.894782 \\
   & R3 & 1.091492 & 0.936947 & 0.862251 & 0.886293 \\
   & R4 & 1.091752 & 0.864618 & 0.864722 & 0.908641 \\
   & R5 & 1.078128 & 0.863782 & 0.917707 & 0.899555 \\
\midrule
P1 & R1 & 1.084886 & 0.901813 & 0.945481 & 0.948446 \\
   & R2 & 1.078187 & 0.904029 & 0.972954 & 0.928288 \\
   & R3 & 1.091492 & 1.001135 & 0.941676 & 0.966558 \\
   & R4 & 1.091752 & 1.029573 & 0.935833 & 0.950611 \\
   & R5 & 1.078128 & 0.970744 & 0.959143 & 0.945202 \\
\midrule
P2 & R1 & 1.199893 & 0.993732 & 0.986854 & 1.008576 \\
   & R2 & 1.192992 & 1.063734 & 1.064368 & 0.989858 \\
   & R3 & 1.206641 & 1.110937 & 0.967441 & 0.996844 \\
   & R4 & 1.206925 & 0.980339 & 0.969911 & 1.026150 \\
   & R5 & 1.192962 & 1.114952 & 1.025218 & 0.992560 \\
\bottomrule
\end{tabular}
\end{table}

在P0的五个输入中，两阶段固定设备联动方法有四题的最终安全能量低于HEFT适配版，其中R2低约10.07\%、R5低约5.88\%，而R3反而高约8.66\%。在P1统一真实反比响应条件下，联动方法仅有R1、R2低于HEFT，R4比HEFT高约10.02\%。当P2的入口温度提高至78\,$^\circ$C后，联动方法仅R2低约0.0595\%，R3则高约14.83\%。因此，三组15个“工作量变体--真实实验条件”组合中，联动低于HEFT的情况为7次，其余8次则未占优。P0/R4的约0.012\%与P2/R2的约0.059\%虽然在高精度区间比较中严格成立，但幅度很小，不能直接推断为工程应用中的稳定优势。

与SA相比，联动方法在15个组合中的8个低于对应五种子的能量中位数。这个数字只说明已观察样本的相对位置，不构成与五次合计1500\,s搜索的等预算比较。联动方法的第二阶段在15组中的14组进一步降低了本次调用第一阶段已找到的完整安全能量，体现了该局部搜索组织在部分输入中能够继续改善已有方案；但不同阶段使用的是同一次方法预算的先后时间，因而该比例不能代替独立、等总预算的阶段消融。

能量质量还需要与求解时间联系起来分析。表\ref{tab:exp-runtime}给出各组五个输入的平均实际方法根时钟。HEFT适配版在15个配对中均比联动MILP更早完成；P0、P1和P2中，联动的平均时间约为HEFT的6.34、6.37和4.99倍。这种差异与两类方法的处理过程相符：HEFT通过优先级、处理器插入和有限局部调整构造完整安排，联动方法还要建立并求解受限的混合整数规划。然而，方法根时钟已包含MATLAB启动、初始化和热评价等开销，不能将全部时间差归因于某一个MILP求解步骤。

\begin{table}[htbp]
\centering
\caption{联动MILP与HEFT的平均方法根时钟（s）}
\label{tab:exp-runtime}
\begin{tabular}{lrrr}
\toprule
实验组 & 联动MILP & HEFT & 平均时钟比 \\
\midrule
P0 & 266.441 & 41.994 & 6.34 \\
P1 & 221.644 & 34.772 & 6.37 \\
P2 & 232.466 & 46.619 & 4.99 \\
\bottomrule
\end{tabular}
\end{table}

个别失利输入还有可定位的搜索范围因素。例如P0/R3中联动方法保留的局部固定配置未包含HEFT已经发现的一份更低频档时间表，P1/R4中第二阶段冻结的部分本地Stage工作点也限制了进一步降档。这些事实能够解释当前局部条件域为什么未直接覆盖相应HEFT方案，却不能证明在所有放宽域条件下都能够取得更优结果。也不能因为整体MILP的初始安全能量较高，就推断完整MILP优化机制本身缺乏节能能力：本批没有取得可评价的整体MILP整数候选，尚不具备这种比较依据。

综合实验一至实验四，本章所建立的模型能够在给定设备和场景条件下区分不同响应信息造成的工作点选择差异、相同执行模式下的时序相关能量变化，以及入口温度对漏电能量和热安全准入的影响。它的实验范围仍有限：P0--P2只覆盖A+B和A+C两种组合、每应用一个请求的三Stage链、四执行单元与5\,ms起点网格，没有对随机DAG拓扑、大规模Stage数量、未知请求在线到达或真实芯片功热参数作系统验证。额外的数据传输和状态切换开销采用零成本假设，且完整能量目标没有纳入规划器运行自身的设备电能。

与此同时，全部120份安全输出并未使完整MILP的独立下界与安全上界闭合，无法给出连续起点联合问题的全局最优证明。正式批次的MILP外层热反馈新增约束行仍为零，也没有反馈带来节能或求解加速的独立证据。已有结果最明确地支持统一电热建模的必要性与连续安全评价的作用；对于在有限计算预算内如何进一步提高安全方案质量、减少搜索成本，需要在相同模型和评价口径下继续研究，这构成第四章的工作重点。

\section{本章小结}

本章面向不同应用组合下的端侧深度学习推理，建立了跨场景共享具体执行单元、场景内分别配置工作点与执行时序的热约束节能联合规划模型。针对CPU、GPU和NPU的合法执行能力及Stage相关频率响应，统一刻画了任务依赖、执行资源互斥和共享频率域约束；进一步结合空闲立即回最低档的运行规则、设备基准和活动功率、本地温度相关漏电以及耦合RC热模型，建立完整评价窗口的系统能量和连续温度安全表达。

为描述这些关系的联合影响，本章将一般连续起点问题与有限启动集合问题加以区分，利用执行选项、固定事件、域状态及温度传播系数形成MILP表达，并通过独立连续热检查判断候选安排的完整安全性。理想模型允许在明确超温时生成有效反馈约束，实际数值求解与安全评价则分别保存其结果和证明范围，避免将有限预算的求解器状态直接等同于完整问题的最优性。

依据已完成的仿真与数学分析，本章观察到任务相关响应信息能够在部分输入中改变频率与时序选择；共享频率域及温度历史使相同Stage模式的合法时间表仍可能具有不同的完整能量；入口温度和温限还会影响漏电能量与候选安全接纳。实验没有观察到MILP外层温度反馈的实际新增约束，也未形成普遍优于HEFT适配版或完整问题全局最优的证据。因此，本章的主要成果是建立并验证一套面向声明端侧抽象平台的统一电热调度评价基础。下一章将在保持任务与电热安全条件一致的前提下，研究如何在有限计算预算内组织更高效的调度搜索，获得能量质量更好的完整安全安排。

\clearpage
\section*{补充推导：固定事件下的热能耗表达}
\label{sec:formal-details}

为说明第3.5节条件传播与积分系数的来源，下面在给定联合工作点和固定片长下展开对应关系。正文的调度选项和频率控制约束负责确定片段的实际状态；本部分推导不新增任务或运行假设。

\subsection*{一、固定控制状态下的物理系数}

在联合工作点状态$\omega$下，各功率组的基准及温度漏电系数由该状态决定。参考温度偏置与节点基准注热向量可写为
\begin{equation}
(b^P_\omega)_n=
\sum_g a^B_{ng}(\xi_g(\omega))B_g(\xi_g(\omega))
+\sum_g\lambda_{gn}(\xi_g(\omega))(T_a-T_g^{\rm ref}).
\label{eq:baseline-vector}
\end{equation}
记$(Q_c)_{no}=h_{n,u_o}q_o$，并令$\ell_{\omega,n}=\sum_g\lambda_{gn}(\xi_g(\omega))$，则节点注热、状态方程系数和片内电功率可整理为
\begin{equation}
\begin{aligned}
p^{\rm heat}&=b^P_\omega+\operatorname{diag}(\ell_\omega)\theta+Q_cv,\\
b_\omega&=C^{-1}b^P_\omega,\qquad D_c=C^{-1}Q_c,\\
\beta_\omega&=\mathbf1^{\mathsf T}b^P_\omega,\\
P^{\rm net}&=\beta_\omega+\ell_\omega^{\mathsf T}\theta+q_c^{\mathsf T}v.
\end{aligned}
\label{eq:physical-coefficients}
\end{equation}
其中$\mathbf1^{\mathsf T}Q_c=q_c^{\mathsf T}$。$A_\omega$的单位为s$^{-1}$，$b_\omega$及$D_c$给出温度变化率，$\beta_\omega$、$q_c$的单位为W，$\ell_\omega$的单位为W/K。上述关系同时对应式\eqref{eq:local-heat-input}的注热和式\eqref{eq:group-electric-power}的计能，防止一项物理功率在两处重复计算。

\subsection*{二、固定片段内的温度传播与能量积分}

当联合状态$\omega$及片段活动向量$v$均不变时，热方程为常系数仿射系统。对长度为$\delta$的片段，定义
\begin{equation}
\begin{aligned}
\Phi_\omega(\delta)&=e^{A_\omega\delta},\\
\mathsf K_\omega(\delta)&=\int_0^\delta e^{A_\omega t}\,dt,\\
\mathsf L_\omega(\delta)&=\int_0^\delta(\delta-t)e^{A_\omega t}\,dt.
\end{aligned}
\label{eq:matrix-integrals}
\end{equation}
温度轨迹因此为
\begin{equation}
\theta(t)=e^{A_\omega t}\theta(0)
+\mathsf K_\omega(t)(b_\omega+D_cv).
\label{eq:segment-solution}
\end{equation}
将温度解代入片内电功率并积分，可得
\begin{equation}
\begin{aligned}
\theta(\delta)&=\Phi_\omega(\delta)\theta(0)
+\mathsf K_\omega(\delta)b_\omega
+\mathsf K_\omega(\delta)D_cv,\\
\Delta E&=\ell_\omega^{\mathsf T}\mathsf K_\omega(\delta)\theta(0)
+\beta_\omega\delta+\ell_\omega^{\mathsf T}\mathsf L_\omega(\delta)b_\omega\\
&\quad+\left[\delta q_c^{\mathsf T}
+\ell_\omega^{\mathsf T}\mathsf L_\omega(\delta)D_c\right]v.
\end{aligned}
\label{eq:segment-integrals}
\end{equation}
与正文式\eqref{eq:thermal-transition}及式\eqref{eq:energy-transition}对照，传播与积分系数分别为
\begin{equation}
\begin{aligned}
\psi&=\mathsf Kb_\omega,&\Gamma&=\mathsf KD_c,\\
\mu^{\mathsf T}&=\ell_\omega^{\mathsf T}\mathsf K,&
\zeta&=\beta_\omega\delta+\ell_\omega^{\mathsf T}\mathsf Lb_\omega,\\
\nu^{\mathsf T}&=\delta q_c^{\mathsf T}+\ell_\omega^{\mathsf T}\mathsf LD_c.
\end{aligned}
\label{eq:integral-coefficients}
\end{equation}
所有系数均在当前$\omega$、$\delta_{c\ell}$下取值。该表达不需要计算$A_\omega^{-1}$，因而不要求该矩阵可逆。阶段活动、等待与最低档空闲尾段都遵循同一传播公式，只是活动向量和联合工作点有所不同。

\subsection*{三、有效变量界与条件线性化}

有限事件、有限工作点和固定窗口保证相应系数存在有限界。采用相容无穷范数，定义
\begin{equation}
a_* =\max_{\omega\in\Omega}\|A_\omega\|_\infty,
\qquad b_{c*}=\max_{\omega\in\Omega}
\left(\|b_\omega\|_\infty+\|D_c\|_\infty\right).
\label{eq:norm-coefficients}
\end{equation}
因为$\|v\|_\infty\le1$，积分比较可给出对整个窗口有效的保守界
\begin{equation}
\|\theta_c(t)\|_\infty\le H_c,
\quad H_c=e^{a_*W_c}
\left(\|T_c^0-T_a\mathbf1\|_\infty+b_{c*}W_c\right).
\label{eq:state-bound}
\end{equation}
该界对任意合法模式和状态切换成立，包括尚未满足温限的时间表。令
\begin{equation}
\overline P_c=\max_\omega
\left(|\beta_\omega|+\|\ell_\omega\|_1H_c+\|q_c\|_1\right),
\quad|\Delta E_{c\ell}|\le\delta_{c\ell}\overline P_c.
\label{eq:energy-bound}
\end{equation}
则温度状态与段能量的有限变量盒可用于构造条件等式的残差界。例如，式\eqref{eq:thermal-transition}与式\eqref{eq:energy-transition}分别可取
\begin{equation}
\begin{aligned}
M^T_n&=H_c+\|\Phi_{n:}\|_1H_c+|\psi_n|+\|\Gamma_{n:}\|_1,\\
M^E&=\delta_{c\ell}\overline P_c+\|\mu\|_1H_c+|\zeta|+\|\nu\|_1.
\end{aligned}
\label{eq:residual-bounds}
\end{equation}
当状态指示变量为1，等式必须成立；为0时，残差界须覆盖所有允许的变量值。式中的大$M$用于精确表达条件关系，不是允许温度超过上限的容差。实际数值建模可收紧这些界，但不能无依据地排除合法的有限域安排。

对于工作点固定的片段，温度有效范围、组功率和节点注热的物理有效性可以统一写成
\begin{equation}
a_{j\omega}^{\mathsf T}\theta+b_{jc\omega}^{\mathsf T}v+d_{j\omega}\le0.
\label{eq:validity-affine}
\end{equation}
在片段起点和终点可以施加对应的必要约束，片段内部仍需连续检查。窗口末点没有正在执行的Stage，因此各控制域均为$L_h$，其功率有效性应按\emph{最低档联合状态}及$v=0$核查，不能沿用旧的“最后一次工作点”或“域状态保持”假设。对任意条件不等式$w=1\Rightarrow a^{\mathsf T}\theta+b^{\mathsf T}v+d\le0$，取$M=\|a\|_1H_c+\|b\|_1+|d|$即可写成对应的有界线性形式。若内点发现明确违规，可用同一状态下的温度解析式构建必要反馈行。

\subsection*{四、选项、状态与真实时间表之间的对应}

给定受限问题$P_{\mathcal G}$的一份合法安排，每个Stage实例都对应一个唯一模式--起点选项，因而能够构造满足任务分配、期限和资源互斥的$x,y$。真实域状态产生$\rho,w$：正在执行的Stage锁定其工作点，没有正在执行的Stage时非最低档支持式\eqref{eq:idle-auto-low}迫使最低档$L_h$。按真实温度轨迹填写片段端点温度和完整段能量，则各条件传播、积分等式均成立。

反过来，若一个整数解满足模式选择、依赖、互斥和域控制约束，便可恢复一张合法时间表。式\eqref{eq:idle-auto-low}与工作阶段锁定式共同确定每个片段的实际控制状态，既不允许空闲片段任意换至高档，也不允许有工作执行期间换档。给定初始温度后，常系数热方程在每段有唯一解，片段端点状态及能量恰由前述积分表达给出。再经完整连续温度与有效域检查合格，便得到$P_{\mathcal G}$的一份可行方案。

反馈行需要保留所有真实安全安排：若候选在某片段内部确有超温，则该时刻的温限是安全安排必须满足的必要条件；未采用对应联合状态的安排不受该行限制。完整选项赋值排除只应针对已经可靠确认违规的完整场景候选。数值未决不能被当作超温或不可行证明。上述结论是模型的数学对应关系，不自动赋予有限计算预算内的求解器结果全局最优证明。


\clearpage
\section*{本章符号索引}
\markright{本章符号索引}

各符号在正文首次出现时解释。下表按已知输入、实际决策与派生变量分类；时间统一换算为秒，功率、能量与温差分别采用W、J和K。

\begin{longtable}{@{}p{5.3cm}p{6.0cm}>{\raggedright\arraybackslash}p{2.5cm}@{}}
\caption{第三章主要符号}\label{tab:symbols}\\
\toprule
符号 & 含义 & 量纲／类别 \\
\midrule
\endfirsthead
\toprule
符号 & 含义 & 量纲／类别 \\
\midrule
\endhead
$\mathcal A,G_a,V_a,E_a$ & 应用集合、各应用的Stage依赖图 & 输入集合 \\
$\mathcal C,\mathcal R_c,\mathcal I_c,\mathcal E_c$ & 场景、请求、Stage实例及依赖 & 输入集合 \\
$\mathcal J,\tau(c,i),r(c,i),\mathcal O_{cr}$ & 固定Stage身份、实例归属、必需输出 & 输入映射 \\
$r_{cr},d^{\rm rel}_{cr},D_{cr},W_c$ & 释放、期限与固定评价窗口 & s \\
$p_c$ & 完整场景的统计权重 & 无量纲 \\
$\mathcal U_C,\mathcal U_G,\mathcal U_N$ & CPU、GPU与NPU执行单元 & 输入集合 \\
$\mathcal H,\gamma(u),\mathcal K_h,L_h$ & 频率域、所属域、工作点、最低档 & 输入集合 \\
$f_h(k),v_h(k)$ & 配套频率与电压 & Hz、V \\
$\mathcal M_{ci},u(m),\kappa(m)$ & 合法模式、具体设备及工作点 & 输入映射 \\
$d^{\rm ref}_{\tau u},f^{\rm ref}_u,\alpha_{\tau u}$ & 参考时长、频率和响应参数 & s、Hz、比例 \\
$d_\tau(m),q_\tau(m)$ & 模式持续时间、活动功率 & s、W \\
$\mathcal B,\pi(u),\xi_{cg}(t)$ & 功率计量组、单元归属与控制状态 & 输入／派生 \\
$B_g,\lambda_{gn},T_g^{\rm ref}$ & 基准、温度漏电系数和参考温度 & W、W/K、$^\circ$C \\
$\mathcal N,C,G,T_a,T_c^0,T_n^{\max}$ & 节点、热容、导热、初温与温限 & J/K、W/K、$^\circ$C \\
$P_{{\rm ext},c},E_{{\rm ext},c}$ & 网络外公共功率及能量 & W、J \\
\midrule
$u_\tau,m_{ci},s_{ci}$ & 共享设备部署、模式、起点 & 决策 \\
$F_{ci},I_{ci},\iota_{ci},k_{ch}(t)$ & 结束、执行区间、活动与域状态 & 派生 \\
$P_{cg},p^{\rm heat}_{cn},\theta_c,E_c,\bar E$ & 组功率、节点注热、温升和能量 & W、K、J \\
\midrule
$\mathcal G_{cim},\mathcal Q_{ci},\mathcal Q_c$ & 有限起点与Stage选项 & 输入／派生 \\
$t_o,f_o,d_o,q_o,u_o,h_o,\kappa_o$ & 选项执行参数与控制要求 & 已知系数 \\
$e_{c\ell},\delta_{c\ell},\chi_{co\ell}$ & 固定事件、片长及活动覆盖 & s、0/1 \\
$x_{\tau u},y_{co}$ & 共享部署、选项选择 & 二元变量 \\
$\rho_{ch\ell k},w_{c\ell\omega}$ & 控制域和联合状态选择 & 二元变量 \\
$\Omega,v_{c\ell}$ & 相容联合状态、活动选项向量 & 派生 \\
$\theta_{c\ell},\Delta E_{c\ell}$ & 片段端点温升与段能量 & K、J \\
$A_\omega,b_\omega,D_c,\beta_\omega,\ell_\omega$ & 热方程及电功率系数 & 派生系数 \\
$\Phi,\psi,\Gamma,\mu,\zeta,\nu$ & 状态传播与积分参数 & 派生系数 \\
$H_c,\overline P_c,M^T,M^E$ & 数值条件等式所用有效界 & 派生界 \\
$\mathcal H^{(j)},U$ & 热必要行集合、安全方案能量上界 & 求解记录 \\
\bottomrule
\end{longtable}

